% «Sahifa dizayni»: spetsifikatsiya va texnik loyiha (Ö Ğ Ş Ç harflarida).
% XeLaTeX yoki LuaLaTeX bilan kompilyatsiya qilinadi.
\documentclass[11pt]{article}
\usepackage{fontspec}
\setmainfont{Noto Sans}
\setmonofont{DejaVu Sans Mono}[Scale=0.88]
\usepackage{newunicodechar}
\newunicodechar{→}{\ensuremath{\rightarrow}}
\newunicodechar{↔}{\ensuremath{\leftrightarrow}}
\newfontfamily\titlefont{texgyreadventor}[Extension=.otf,UprightFont=*-bold]
\usepackage[uzbek,provide=*]{babel}
\setlocalecaption{uzbek}{contents}{Mundarija}
\setlocalecaption{uzbek}{figure}{Rasm}
\setlocalecaption{uzbek}{table}{Jadval}
\usepackage[a4paper,top=22mm,bottom=22mm,left=20mm,right=20mm,headheight=14pt]{geometry}
\usepackage{microtype}
\usepackage{xcolor}
\usepackage[most]{tcolorbox}
\usepackage{enumitem}
\usepackage{array}
\usepackage{tabularx}
\usepackage{multicol}
\usepackage{booktabs}
\usepackage{caption}
\usepackage{fancyhdr}
\usepackage{titlesec}
\usepackage{hyperref}
\definecolor{asosiy}{HTML}{2A4D6E}
\definecolor{urgu}{HTML}{D9822B}
\definecolor{yashil}{HTML}{2E8B57}
\definecolor{fon}{HTML}{EEF3F8}
\definecolor{kulrang}{HTML}{6B7280}
\definecolor{qizil}{HTML}{B42318}
\hypersetup{colorlinks,linkcolor=asosiy,urlcolor=asosiy}
\setlist{itemsep=1pt,topsep=2pt,leftmargin=16pt}
\setlist[itemize]{label={\color{urgu}\small\textbullet}}
\setlist[enumerate]{label={\color{asosiy}\bfseries\arabic*.}}
\titleformat{\section}{\titlefont\Large\color{asosiy}}{\thesection.}{0.5em}{}[{\color{urgu}\titlerule[1.2pt]}]
\titleformat{\subsection}{\titlefont\large\color{asosiy}}{\thesubsection.}{0.5em}{}
\titlespacing*{\section}{0pt}{8pt}{5pt}
\titlespacing*{\subsection}{0pt}{7pt}{3pt}
\pagestyle{fancy}\fancyhf{}
\fancyhead[L]{\small\color{asosiy}Oversells: Sahifa dizayni}
\fancyhead[R]{\small\color{asosiy}\leftmark}
\fancyfoot[C]{\small\thepage}
\renewcommand\headrulewidth{0.4pt}
\renewcommand\sectionmark[1]{\markboth{#1}{}}
\setlength{\parindent}{0pt}\setlength{\parskip}{2pt}\setlength{\emergencystretch}{5em}
\DeclareCaptionFont{asosiyrang}{\color{asosiy}}
\captionsetup{font=small,labelfont={bf,asosiyrang},labelsep=period,skip=3pt}
\setlength{\textfloatsep}{8pt plus 2pt}\setlength{\floatsep}{8pt plus 2pt}\setlength{\intextsep}{8pt plus 2pt}
\renewcommand\arraystretch{1.15}
\setcounter{tocdepth}{1}
% muharrirdagi tugma nomi
\newcommand\tugma[1]{\tcbox[on line,boxsep=0pt,left=3pt,right=3pt,top=1.5pt,bottom=1.5pt,colback=fon,colframe=asosiy!50,boxrule=0.5pt,arc=2pt]{\small\textbf{#1}}}
% maslahat qutisi
\newtcolorbox{maslahat}[1][Maslahat]{colback=urgu!8,colframe=urgu,boxrule=0.8pt,arc=2mm,left=3mm,right=3mm,top=1.5mm,bottom=1.5mm,fonttitle=\bfseries,title=#1}
% ---------------------------------------------------------------------
% Holat belgilari. Har bir imkoniyat yonida bittasi turadi:
%   \tayyor  --- işlaydi;
%   \beta    --- sinovda (özgariși mumkin);
%   \reja    --- hali yöq, rejada.
% Holatni almaştiriş: imkoniyat qatoridagi buyruqni boşqasiga almaştiring
% (masalan, \tayyor örniga \beta). Belgining rangi va yozuvi faqat
% \holatbelgi da özgartiriladi.
% \beta yunon harfi örniga işlatiladi; formulada beta kerak bölsa,
% \betaharf ni yozing.
\newcommand\holatbelgi[2]{\tcbox[on line,boxsep=0pt,left=2.5pt,right=2.5pt,top=1pt,bottom=1pt,colback=#1!12,colframe=#1,boxrule=0.4pt,arc=2pt]{\scriptsize\textbf{\color{#1!75!black}#2}}}
\newcommand\tayyor{\holatbelgi{yashil}{tayyor}}
\let\betaharf\beta
\renewcommand\beta{\holatbelgi{urgu}{beta}}
\newcommand\reja{\holatbelgi{kulrang}{reja}}
% tekşirilişi kerak bölgan fakt
\newcommand\tekshir[1]{\leavevmode{\color{qizil}\small[TEKSHIRILSIN: #1]}}
% rasm uçun joy: #1 balandlik, #2 ekranda nima körinişi, #3 sarlavha
\newcommand\rasmjoy[3]{%
  \begin{figure}[!htbp]\centering\vspace{-2pt}
  \fbox{\begin{minipage}[c][#1][c]{0.9\linewidth}\centering\small\color{asosiy}\textbf{Rasm uçun joy}\par\smallskip #2\end{minipage}}
  \caption{#3}
  \end{figure}}

\usepackage{fancyvrb}
\fvset{fontsize=\footnotesize,frame=single,rulecolor=\color{asosiy!40},framesep=2mm}

\begin{document}

\begin{tcolorbox}[colback=asosiy,colframe=asosiy,arc=3mm,left=6mm,right=6mm,top=4mm,bottom=4mm]
\color{white}
{\titlefont\fontsize{26}{30}\selectfont Sahifa dizayni}\par\medskip
{\large Har qanday kitobning içki sahifalari: spetsifikatsiya, texnik loyiha va yöl xaritasi}
\end{tcolorbox}

Bu hujjat Oversells'ga qöşiladigan «Sahifa dizayni» imkoniyatlarini tasvirlaydi. A qismi (\ref{sec:sd}--\ref{sec:sd-ai}-bölimlar) imkoniyatlar hujjatiga özgarişsiz qöşiladi: u \texttt{bolimlar.tex} faylida. Keyingi bölimlar --- texnik loyiha, foydalanuvçi yöllari, xavflar, yöl xaritasi va savollar.
\textbf{Belgilar:} \tayyor\ --- işlaydi; \beta\ --- sinovda; \reja\ --- rejada. Qizil \tekshir{...} --- tasdiqlanmagan maʼlumot.

\begin{multicols}{2}\small
\tableofcontents
\end{multicols}

\part*{A. Spetsifikatsiya}
% =====================================================================
%  «Sahifa dizayni» --- imkoniyatlar hujjatiga qöşiladigan bölimlar.
%  docs/imkoniyatlar/oversells-imkoniyatlari.tex dagi «Bosmaxona va
%  eksport» bölimidan keyin % =====================================================================
%  «Sahifa dizayni» --- imkoniyatlar hujjatiga qöşiladigan bölimlar.
%  docs/imkoniyatlar/oversells-imkoniyatlari.tex dagi «Bosmaxona va
%  eksport» bölimidan keyin % =====================================================================
%  «Sahifa dizayni» --- imkoniyatlar hujjatiga qöşiladigan bölimlar.
%  docs/imkoniyatlar/oversells-imkoniyatlari.tex dagi «Bosmaxona va
%  eksport» bölimidan keyin \input{../sahifa-dizayni/bolimlar} bilan
%  qöşiladi. Preambuladagi \tugma, \maslahat, \tayyor, \beta, \reja,
%  \tekshir va \rasmjoy işlatiladi. Hamma imkoniyat hozir \reja.
% =====================================================================

\section{Sahifa dizayni: umumiy körinişi}\label{sec:sd}
\begin{itemize}
\item \textbf{Nima} \reja\ --- kitob içki sahifalarini siçqonça bilan joylaş: hoşiyalar, ustunlar, rasm ramkalari, matn qutilari, sahifa ustalari va uslublar. Har bir özgariş LaTeX kodiga yoziladi.
\item \textbf{Qanday kitoblar} \beta\ --- roman va nasr, şeʼriyat, bolalar rasmli kitobi, luğat, diniy matn (arabça, öngdan çapga), oşpazlik kitobi, iş daftari, fotoalbom, katalog.
\item \textbf{Natija} \reja\ --- tahrirlanadigan LaTeX va bosmaga tayyor PDF. Yopiq yoki ikkilik format yöq: barça sozlamalar loyihadagi \texttt{.sty} va \texttt{.tex} fayllarda turadi.
\item \textbf{Özbek tili} \reja\ --- yangi alifbo (Ö Ğ Ş Ç) uçun imlo, özbekça böğin köçiriş va kirill ↔ lotin ötkaziş dizayn qilingan sahifalarda ham işlaydi.
\item \textbf{Kompilyator} \beta\ --- XeLaTeX va LuaLaTeX; arabça (öngdan çapga) şablon ikkalasida ham işlaydi, LuaLaTeX tavsiya etiladi (\ref{sec:sd-shablon}-bölim).
\end{itemize}

\rasmjoy{1.2cm}{Ekran: çapda kod, öngda kompilyatsiya qilingan sahifa; sahifa ustida kök hoşiya çiziqlari, ustun çegaralari va rasm ramkasining tutqiçlari; yuqorida «Sahifa», «Ustalar», «Uslublar», «Tipografiya», «Preflight» yorliqlari.}{Sahifa dizayni rejimi}

% ---------------------------------------------------------------------
\section{Körinadigan sahifa rejimi}\label{sec:sd-vizual}
\begin{itemize}
\item \textbf{\tugma{Sahifa dizayni}} \reja\ --- PDF panelini tahrir rejimiga ötkazadi: kod çapda, sahifa öngda qoladi.
\item \textbf{Hoşiya va ustunlar} \beta\ --- hoşiya çiziqlari va ustun çegaralari siçqonça bilan suriladi; qiymat millimetrda yonida körinadi.
\item \textbf{Ramkalar} \beta\ --- rasm ramkasi va matn qutisi suriladi, kattalaştiriladi, kesiladi.
\item \textbf{Koddan sahifaga} \beta\ --- kodni özgartirsangiz, kompilyatsiyadan keyin sahifa va tutqiçlar yangilanadi.
\item \textbf{Sahifadan kodga} \beta\ --- tutqiçni qöyib yuborganingizda faqat şu ramkaning kalitlari (\texttt{x=}, \texttt{y=}, \texttt{w=}\ldots) özgaradi; qolgan kod harfma-harf joyida qoladi.
\item \textbf{Siz yozgan kod} \beta\ --- vizual tahrirlovçiga tegişli bölmagan kod sahifada qulf belgisi bilan körinadi; ustiga bossangiz, kodning öşa qatori oçiladi.
\item \textbf{Birgalikda} \beta\ --- özgarişlar boşqa tahrirlar kabi jonli sinxron qilinadi va «Tarix» ga tuşadi.
\end{itemize}

\textbf{Vizual tahrirlovçiga tegişli LaTeX qurilmalari}
\begin{center}\footnotesize
\begin{tabularx}{\linewidth}{>{\raggedright\arraybackslash}p{50mm}X}
\toprule
Qurilma & Nimani boşqaradi\\
\midrule
\texttt{sahifa.sty} (butunlay) & sahifa ölçami, hoşiyalar, ustunlar, ustalar, uslublar, şriftlar; boşida \texttt{\% oversells-sahifa:} izohida model\\
\texttt{\textbackslash sahifausta\{id\}} & şu joydan boşlab qaysi usta işlaşi\\
\texttt{\textbackslash begin\{uslub\}\{id\}} & xatboşi uslubi\\
\texttt{\textbackslash matnuslubi\{id\}\{...\}} & belgi uslubi\\
\texttt{\textbackslash rasmramka[...]\{fayl\}} & rasm: oqimda, matn bilan oralaş, sahifada joylaşgan yoki töliq sahifa\\
\texttt{\textbackslash begin\{matnquti\}[...]} & sahifaga mahkamlangan matn qutisi\\
\texttt{\textbackslash begin\{ustunlar\}[...]\{n\}} & bir qismdagi ustunlar soni\\
\texttt{\textbackslash bosbet} & ataylab böş qoldirilgan bet\\
\bottomrule
\end{tabularx}
\end{center}
Boşqa hamma narsa (sizning buyruqlaringiz, \texttt{tikzpicture}, jadvallar) faqat kod orqali tahrirlanadi.

% ---------------------------------------------------------------------
\section{Sahifa ustalari}\label{sec:sd-usta}
\begin{itemize}
\item \textbf{Juft va toq betlar} \beta\ --- har biriga alohida yuqori va pastki yozuvlar: bob nomi, bölim nomi, kitob nomi, bet raqami yoki matn.
\item \textbf{Bet raqami} \beta\ --- joyi (yuqori/past, içki/taşqi çet, örta), şakli (1, i, I), boşlaniş raqami.
\item \textbf{Bob boşlanişi} \beta\ --- bob doim toq betda boşlanadimi; bob betida yuqori yozuv körinadimi; bob raqami va bezagi.
\item \textbf{Böş bet qoidasi} \beta\ --- bob oldidan qolgan böş betda yuqori va pastki yozuvlar bölmaydi.
\item \textbf{Fon elementlari} \beta\ --- ustaga rasm, rangli maydon yoki çiziq qöyiş; şu usta işlagan har bir betda çiqadi, juft betlarda kerak bölsa közgu kabi aks etadi; bob oldidagi böş betlarda çiqmaydi.
\item \textbf{Bölimga xos usta} \beta\ --- muqaddima, asosiy matn, ilovalar uçun turli ustalar; \texttt{\textbackslash sahifausta\{id\}} bilan almaşadi.
\item \textbf{Butun kitobga} \beta\ --- ustani bir marta özgartirsangiz, u işlatilgan hamma bet yangilanadi.
\end{itemize}

\rasmjoy{1.2cm}{Ekran: «Ustalar» paneli --- ikki betli yoyilma (juft/toq), yuqori yozuv kataklari: çapda «bob nomi», öngda «bet raqami»; röyxatda «Asosiy», «Bob boşlanişi», «Ilova» ustalari.}{Sahifa ustalari}

% ---------------------------------------------------------------------
\section{Uslublar}\label{sec:sd-uslub}
\begin{itemize}
\item \textbf{Xatboşi uslublari} \beta\ --- Sarlavha 1--3, Asosiy matn, Şeʼr qatori, Iqtibos, Izoh, Rasm osti yozuvi va öz uslublaringiz.
\item \textbf{Belgi uslublari} \beta\ --- urğu, atama, xorijiy söz, sahifa havolasi va öz uslublaringiz.
\item \textbf{Qöllaş} \beta\ --- matnni belgilab uslubni tanlaysiz; uslub LaTeX buyruği bölib kodga yoziladi.
\item \textbf{Bir joyda özgartiriş} \beta\ --- uslubni özgartirsangiz, kitobdagi hamma joy yangilanadi, çunki uslub \texttt{sahifa.sty} da bir marta taʼriflangan.
\item \textbf{Asosiy uslub} \beta\ --- uslub boşqasidan meros oladi («Iqtibos» = «Asosiy matn» + kursiv + surilma).
\end{itemize}

\begin{center}\footnotesize
\begin{tabularx}{\linewidth}{>{\bfseries}p{30mm}>{\raggedright\arraybackslash}p{58mm}X}
\toprule
Uslub & LaTeX'da & Izoh\\
\midrule
Sarlavha 1, 2, 3 & \texttt{\textbackslash chapter}, \texttt{\textbackslash section}, \texttt{\textbackslash subsection} & körinişi titlesec bilan; mundarija, «Faqat şu bob» va SyncTeX özgarmaydi\\
Asosiy matn & hujjatning odatiy şrifti & \texttt{\textbackslash setmainfont}, \texttt{\textbackslash fontsize}\\
Şeʼr qatori & \texttt{\textbackslash begin\{uslub\}\{sher\}} & qatorlar \texttt{\textbackslash\textbackslash} bilan; uzun qator surilib davom etadi\\
Iqtibos, Izoh & \texttt{\textbackslash begin\{uslub\}\{iqtibos\}} & ixtiyoriy nomli uslub ham şu yöl bilan\\
Rasm osti yozuvi & caption paketi sozlamasi & \texttt{\textbackslash rasmramka[izoh=...]}\\
Belgi uslubi & \texttt{\textbackslash matnuslubi\{urgu\}\{...\}} & \\
\bottomrule
\end{tabularx}
\end{center}

% ---------------------------------------------------------------------
\section{Tipografiya paneli}\label{sec:sd-tipo}
\begin{center}\footnotesize
\begin{tabularx}{\linewidth}{>{\bfseries}p{34mm}X>{\raggedright\arraybackslash}p{40mm}}
\toprule
Sozlama & Nima qiladi & LaTeX'da\\
\midrule
Şrift & yuklangan TTF/OTF yoki örnatilgan şrift, har bir uslubga alohida & \texttt{\textbackslash newfontfamily} (fontspec)\\
Ölçam va qator oralişi & punkt bilan & \texttt{\textbackslash fontsize\{..\}\{..\}}\\
Harf oralişi (tracking) & uslub uçun & fontspec \texttt{LetterSpace}\\
Tayanç çiziq töri & vizual rejimda tör çiziqlari körinadi (\beta); joylarni avtomatik qator oralişiga karrali qiliş \reja & \ref{sec:sd-xavf}-bölimdagi çeklov bilan\\
Böğin köçiriş & yoqiş/öçiriş, eng kam harf soni; özbekça qoidalar (lotin, kirill, Ö Ğ Ş Ç) & \texttt{\textbackslash uzbekhyphenmins}, \texttt{\textbackslash hyphenpenalty}\\
Yolğiz qatorlar & bet boşida yoki oxirida bitta qator qolmasin & \texttt{\textbackslash widowpenalty}, \texttt{\textbackslash clubpenalty}; LuaLaTeX'da lua-widow-control\\
Bosh harf (drop cap) & 2--4 qator balandlikda & lettrine\\
Bob bezaklari & bob sarlavhasi ostida va bob oxirida \texttt{* * *}, qisqa çiziq yoki gul (\beta) & \texttt{\textbackslash bobbezagi}, pifont\\
Mikrotipografiya & çetlarni tekislaş, harflarni ozgina keng/tor qiliş & microtype (XeLaTeX'da faqat çetni tekislaş)\\
\bottomrule
\end{tabularx}
\end{center}

% ---------------------------------------------------------------------
\section{Rasm ramkalari va matn qutilari}\label{sec:sd-ramka}
\begin{itemize}
\item \textbf{Taşlab qöyiş} \reja\ --- rasmni sahifaga taşlaysiz; \texttt{\textbackslash rasmramka} yoziladi, rasm \texttt{rasmlar/} papkasiga tuşadi (hozirgi rasm qöyiş qoidalari bilan).
\item \textbf{Kesiş} \reja\ --- ramka çetlarini Alt bilan tortib kesiladi; graphicx \texttt{trim} va \texttt{clip} ga yoziladi, asl fayl özgarmaydi.
\item \textbf{Matn bilan oralaş} \beta\ --- rasm çap yoki öngda, matn uni aylanib ötadi (wrapfig).
\item \textbf{Sahifaga mahkamlaş} \beta\ --- rasm yoki matn qutisi betning aniq joyida turadi (textpos, eso-pic).
\item \textbf{Töliq bet va çetgaça} \beta\ --- rasm butun betni egallaydi va kesiş çizigidan 3 mm taşqariga çiqadi (bleed).
\item \textbf{Rasm osti yozuvi} \beta\ --- ramka ostida yoki yonida; uslubi «Rasm osti yozuvi».
\item \textbf{Rasmli kitob maketi} \beta\ --- töliq betli rasm (\texttt{joy=fon}) + ustida matn qutisi (\texttt{matnquti}); ikki betli yoyilmaga bitta rasm (\texttt{joy=yoyilma}).
\end{itemize}

\textbf{Rasm ramkasi rejimlari}
\begin{center}\footnotesize
\begin{tabularx}{\linewidth}{>{\bfseries}p{28mm}X>{\raggedright\arraybackslash}p{36mm}}
\toprule
Rejim & Qanday turadi & Paket\\
\midrule
\texttt{joy=oqim} & matn içida, xatboşilar orasida & graphicx\\
\texttt{joy=oralash} & çap yoki öngda, matn aylanib ötadi & wrapfig\\
\texttt{joy=sahifa} & betning aniq koordinatasida & textpos, eso-pic\\
\texttt{joy=tola} & butun bet, çetgaça; bet raqami va yuqori yozuv yöq & eso-pic\\
\texttt{joy=yoyilma} & ikki betga bitta rasm (çap yarmi juft betda, öng yarmi toq betda) & eso-pic, graphicx \texttt{trim}\\
\bottomrule
\end{tabularx}
\end{center}

\rasmjoy{1.2cm}{Ekran: bolalar kitobi yoyilmasi --- ikki betni egallagan rasm, kesiş çizigidan taşqariga çiqqan qizil bleed hoşiyasi, rasm ustida yarim şaffof fonli matn qutisi va uning tutqiçlari.}{Töliq betli rasm va matn qutisi}

% ---------------------------------------------------------------------
\section{Kitob turlari uçun şablonlar}\label{sec:sd-shablon}
Har bir şablon özbekça (lotin, kirill, Ö Ğ Ş Ç) va xatosiz kompilyatsiya qilinişi şart. \reja
\begin{center}\footnotesize
\begin{tabularx}{\linewidth}{>{\bfseries}p{28mm}X>{\raggedright\arraybackslash}p{30mm}}
\toprule
Şablon & Nima bor & Kompilyator\\
\midrule
Roman, nasr \beta & bir ustun, bob boşida bosh harf, bob oxirida bezak, toq betda boşlanadigan bob & XeLaTeX, LuaLaTeX\\
Şeʼriyat \beta & şeʼr qatori uslubi, raqamsiz betlar, şeʼr markazda, uzun qator surilib davom etadi & XeLaTeX, LuaLaTeX\\
Bolalar rasmli kitobi \beta & töliq betli rasmlar, katta şrift, rasm ustida matn qutisi, kvadrat format & XeLaTeX, LuaLaTeX\\
Luğat \beta & ikki ustun, qalin bosh söz, betning yuqorisida birinçi va oxirgi bosh söz, harf sarlavhalari & XeLaTeX, LuaLaTeX\\
Diniy matn \beta & arabça matn öngdan çapga, özbekça tarjima yonma-yon (paracol), oyat raqami & LuaLaTeX, XeLaTeX\\
Oşpazlik kitobi \beta & «Masalliqlar» va «Tayyorlaş» ikki ustunda, taom rasmi & XeLaTeX, LuaLaTeX\\
Iş daftari \beta & yoziş çiziqlari, javob uçun joy, katak qoğoz (mavjud dizayn paketi bilan) & XeLaTeX, LuaLaTeX\\
Fotoalbom, katalog \beta & rasm töri (2, 3, 4 rasm; \texttt{rasmtori}, \texttt{\textbackslash torrasm}), rasm osti yozuvi, mahsulot kartoçkasi & XeLaTeX, LuaLaTeX\\
\bottomrule
\end{tabularx}
\end{center}

% ---------------------------------------------------------------------
\section{Muqova dizayneri}\label{sec:sd-muqova}
\begin{itemize}
\item \textbf{Üç qism} \beta\ --- old muqova, orqa muqova va kitob bel qismi (koreşok) bitta yoyilmada.
\item \textbf{Bel kengligi} \beta\ --- bet soni va qoğoz qalinligidan hisoblanadi: \(\text{kenglik} = \frac{\text{betlar}}{2} \times \text{bir varaq qalinligi}\) + muqova qöşimçasi.
\item \textbf{Bleed va kesiş çiziqlari} \beta\ --- har tomonda bleed (masalan 3 mm) va kesiş belgilari.
\item \textbf{CMYK va ISBN} \beta\ --- ranglar CMYK; ISBN ştrix-kodi «Bosmaxona» modulidagi EAN-13 generatori bilan.
\item \textbf{Alohida PDF} \beta\ --- muqova \texttt{muqova.tex} dan alohida kompilyatsiya qilinadi va alohida yuklab olinadi.
\end{itemize}

\textbf{Qoğoz qalinligi jadvali} \beta\ --- gazeta (45, 49), ofset (60--120), bör qoğoz yaltiroq va xira (90--170), kitob qoğozi (60--80 g/m²) uçun odatiy qalinliklar; ular taxminiy (zavodga qarab $\pm$10\,\%), şuning uçun bosmaxona bergan aniq qiymat alohida kiritiladi. Muqovalaş: yelimlangan yumşoq, skoba, qattiq (karton 2,5 mm, 3 mm çiqib turadi, 7 mm bukilma, 15 mm qayirma --- odatiy qiymatlar). \tekshir{qattiq muqova ölçamlarini bosmaxona bilan solişriş}

\rasmjoy{1.2cm}{Ekran: muqova yoyilmasi --- çapda orqa muqova (annotatsiya, ştrix-kod), örtada 14 mm bel (kitob nomi tik yozilgan), öngda old muqova; atrofida bleed va kesiş belgilari; yonida «Bet soni: 256, qoğoz: 80 g/m²» maydonlari.}{Muqova dizayneri}

% ---------------------------------------------------------------------
\section{Preflight: bosmadan oldingi tekşiruv}\label{sec:sd-preflight}
\tugma{Preflight} kompilyatsiya qilingan PDF ni tekşiradi va röyxat körsatadi; tuzatib bölsa, yonida tugma. \beta
\begin{center}\footnotesize
\begin{tabularx}{\linewidth}{>{\bfseries}p{36mm}X>{\raggedright\arraybackslash}p{40mm}}
\toprule
Tekşiruv & Qanday aniqlanadi & Bir bosişda tuzatiş\\
\midrule
Şriftlar joylanganmi & \texttt{pdffonts} & yöq (XeLaTeX/LuaLaTeX şriftlarni odatda joylaydi)\\
Rasm sifati (dpi) & \texttt{pdfimages -list}: joylangan ölçamdagi piksel ziçligi & yöq; rasm nomi va beti körsatiladi\\
Rasm rangi (RGB/CMYK) & \texttt{pdfimages -list} & yöq; bosmaxonaga ögiriş kerakligi aytiladi\\
Bleed & log da çetgaça rasm (\texttt{tola}, \texttt{fon}, \texttt{yoyilma}) bor, lekin PDF beti kesilgan betdan katta emas & yöq; «Bosmaxona»da kesiş belgilarini yoqiş aytiladi\\
Bet soni 16 yoki 32 ga karrali & log dagi bet soni & oxiriga kerakli böş betlar qöşiş\\
Yolğiz qatorlar va qisqa betlar & yolğiz qatorlarni jarimalar oldini oladi; bet oxirida 1--2 qator qolgan betlar \texttt{pdftotext -layout} qatorlari bilan topiladi & yöq; \texttt{\textbackslash enlargethispage} maslahati\\
Sahifadan çiqqan qatorlar & log dagi «Overfull \textbackslash hbox» & yöq; qatorga ötiş\\
Böş betlar & matni ham, rasmi ham yöq bet (\texttt{\textbackslash bosbet} dan taşqari) & yöq; betga ötiş\\
\bottomrule
\end{tabularx}
\end{center}
Hammasi işlaydi: şriftlar, rasm sifati va rangi, bet soni (böş betlarni qöşiş tugmasi bilan), sahifadan çiqqan qatorlar, böş betlar, qisqa betlar va bleed.

% ---------------------------------------------------------------------
\section{Import}\label{sec:sd-import}
\begin{itemize}
\item \textbf{Word (.docx)} \beta\ --- matn, rasmlar va formulalar hozirgi «Wordʼdan import» kabi; qöşimça ravişda Word xatboşi uslublari (Heading 1, Quote, öz uslublaringiz) Oversells uslublariga aylanadi.
\item \textbf{InDesign (.idml)} \beta\ --- ramkalar sahifadagi tartibida öqiladi; matn, xatboşi va belgi uslublari (ölçam, qator oralişi, tekislaş, surilmalar, oraliqlar, qalin va kursiv, rang, bosh harf, meros), sarlavhalar, izohlar, sahifa ölçami va birinçi ustaning hoşiyalari ötadi. Ötmaydi: ramkalarning aniq joyi, jadval (matni xatboşi bölib qoladi), boğlangan rasmlar va şriftlar --- ular röyxati \texttt{main.tex} boşida.
\item \textbf{PDF} \beta\ --- faqat matn (\texttt{pdftotext -layout} sandbox içida): xatboşilar surilma va böş qatordan topiladi, takrorlanadigan sarlavha va bet raqamlari olib taşlanadi, satr oxiridagi bölinişlar tiklanadi («bir-biri» kabi juft sözlar saqlanadi), sarlavhalar taxminiy (\texttt{\textbackslash section*}). Uslublar, rasmlar va jadvallar ötmaydi; skanerlangan PDF da matn yöq.
\item \textbf{Uslublar moslaşuvi} \beta\ --- «Uslublar» bölimida istalgan uslubni boşqasiga qöşasiz (masalan, Word'dan kelgan «Quote» → «Iqtibos»): loyihadagi barça \texttt{.tex} fayllarda u işlatilgan joylar özgaradi, uslub \texttt{sahifa.sty} dan öçiriladi, unga asoslangan uslublar yangisiga asoslanadi.
\end{itemize}

% ---------------------------------------------------------------------
\section{AI dizayn}\label{sec:sd-ai}
Faqat server egasi AI ni yoqqan bölsa işlaydi; yoqilmagan bölsa, tugmalar körinmaydi va heç narsa taşqariga yuborilmaydi. \beta
\begin{itemize}
\item \textbf{Kitobdan dizayn} \reja\ --- mavjud kitobning 2--3 betini (PDF yoki rasm) yuklaysiz; uslublar, sahifa ölçami, hoşiyalar va ustalar taklif qilinadi. Hozirgi «Kitobdan dizayn» imkoniyatining kengaytmasi.
\item \textbf{Tavsifdan dizayn} \beta\ --- «bolalar kitobi, iliq ranglar, katta şrift» kabi tavsif yozasiz; dizayn taklif qilinadi.
\item \textbf{Taklif --- faqat taklif} \beta\ --- natija körinadi, qabul qilsangizgina \texttt{sahifa.sty} ga yoziladi; «Tarix» da oldingi holat saqlanadi.
\end{itemize}
\begin{center}\footnotesize
\begin{tabularx}{\linewidth}{>{\bfseries}p{40mm}p{40mm}X}
\toprule
Imkoniyat & Qayerda & Taşqariga nima ketadi\\
\midrule
AI: kitobdan sahifa dizayni & taşqi AI serveri (Anthropic) & siz yuklagan betlar (PDF yoki rasm), istaklaringiz, şriftlar röyxati\\
AI: tavsifdan dizayn & taşqi AI serveri & siz yozgan tavsif va şriftlar röyxati\\
\bottomrule
\end{tabularx}
\end{center}
Kitob matni, loyiha fayllari va boşqa betlar yuborilmaydi.
 bilan
%  qöşiladi. Preambuladagi \tugma, \maslahat, \tayyor, \beta, \reja,
%  \tekshir va \rasmjoy işlatiladi. Hamma imkoniyat hozir \reja.
% =====================================================================

\section{Sahifa dizayni: umumiy körinişi}\label{sec:sd}
\begin{itemize}
\item \textbf{Nima} \reja\ --- kitob içki sahifalarini siçqonça bilan joylaş: hoşiyalar, ustunlar, rasm ramkalari, matn qutilari, sahifa ustalari va uslublar. Har bir özgariş LaTeX kodiga yoziladi.
\item \textbf{Qanday kitoblar} \beta\ --- roman va nasr, şeʼriyat, bolalar rasmli kitobi, luğat, diniy matn (arabça, öngdan çapga), oşpazlik kitobi, iş daftari, fotoalbom, katalog.
\item \textbf{Natija} \reja\ --- tahrirlanadigan LaTeX va bosmaga tayyor PDF. Yopiq yoki ikkilik format yöq: barça sozlamalar loyihadagi \texttt{.sty} va \texttt{.tex} fayllarda turadi.
\item \textbf{Özbek tili} \reja\ --- yangi alifbo (Ö Ğ Ş Ç) uçun imlo, özbekça böğin köçiriş va kirill ↔ lotin ötkaziş dizayn qilingan sahifalarda ham işlaydi.
\item \textbf{Kompilyator} \beta\ --- XeLaTeX va LuaLaTeX; arabça (öngdan çapga) şablon ikkalasida ham işlaydi, LuaLaTeX tavsiya etiladi (\ref{sec:sd-shablon}-bölim).
\end{itemize}

\rasmjoy{1.2cm}{Ekran: çapda kod, öngda kompilyatsiya qilingan sahifa; sahifa ustida kök hoşiya çiziqlari, ustun çegaralari va rasm ramkasining tutqiçlari; yuqorida «Sahifa», «Ustalar», «Uslublar», «Tipografiya», «Preflight» yorliqlari.}{Sahifa dizayni rejimi}

% ---------------------------------------------------------------------
\section{Körinadigan sahifa rejimi}\label{sec:sd-vizual}
\begin{itemize}
\item \textbf{\tugma{Sahifa dizayni}} \reja\ --- PDF panelini tahrir rejimiga ötkazadi: kod çapda, sahifa öngda qoladi.
\item \textbf{Hoşiya va ustunlar} \beta\ --- hoşiya çiziqlari va ustun çegaralari siçqonça bilan suriladi; qiymat millimetrda yonida körinadi.
\item \textbf{Ramkalar} \beta\ --- rasm ramkasi va matn qutisi suriladi, kattalaştiriladi, kesiladi.
\item \textbf{Koddan sahifaga} \beta\ --- kodni özgartirsangiz, kompilyatsiyadan keyin sahifa va tutqiçlar yangilanadi.
\item \textbf{Sahifadan kodga} \beta\ --- tutqiçni qöyib yuborganingizda faqat şu ramkaning kalitlari (\texttt{x=}, \texttt{y=}, \texttt{w=}\ldots) özgaradi; qolgan kod harfma-harf joyida qoladi.
\item \textbf{Siz yozgan kod} \beta\ --- vizual tahrirlovçiga tegişli bölmagan kod sahifada qulf belgisi bilan körinadi; ustiga bossangiz, kodning öşa qatori oçiladi.
\item \textbf{Birgalikda} \beta\ --- özgarişlar boşqa tahrirlar kabi jonli sinxron qilinadi va «Tarix» ga tuşadi.
\end{itemize}

\textbf{Vizual tahrirlovçiga tegişli LaTeX qurilmalari}
\begin{center}\footnotesize
\begin{tabularx}{\linewidth}{>{\raggedright\arraybackslash}p{50mm}X}
\toprule
Qurilma & Nimani boşqaradi\\
\midrule
\texttt{sahifa.sty} (butunlay) & sahifa ölçami, hoşiyalar, ustunlar, ustalar, uslublar, şriftlar; boşida \texttt{\% oversells-sahifa:} izohida model\\
\texttt{\textbackslash sahifausta\{id\}} & şu joydan boşlab qaysi usta işlaşi\\
\texttt{\textbackslash begin\{uslub\}\{id\}} & xatboşi uslubi\\
\texttt{\textbackslash matnuslubi\{id\}\{...\}} & belgi uslubi\\
\texttt{\textbackslash rasmramka[...]\{fayl\}} & rasm: oqimda, matn bilan oralaş, sahifada joylaşgan yoki töliq sahifa\\
\texttt{\textbackslash begin\{matnquti\}[...]} & sahifaga mahkamlangan matn qutisi\\
\texttt{\textbackslash begin\{ustunlar\}[...]\{n\}} & bir qismdagi ustunlar soni\\
\texttt{\textbackslash bosbet} & ataylab böş qoldirilgan bet\\
\bottomrule
\end{tabularx}
\end{center}
Boşqa hamma narsa (sizning buyruqlaringiz, \texttt{tikzpicture}, jadvallar) faqat kod orqali tahrirlanadi.

% ---------------------------------------------------------------------
\section{Sahifa ustalari}\label{sec:sd-usta}
\begin{itemize}
\item \textbf{Juft va toq betlar} \beta\ --- har biriga alohida yuqori va pastki yozuvlar: bob nomi, bölim nomi, kitob nomi, bet raqami yoki matn.
\item \textbf{Bet raqami} \beta\ --- joyi (yuqori/past, içki/taşqi çet, örta), şakli (1, i, I), boşlaniş raqami.
\item \textbf{Bob boşlanişi} \beta\ --- bob doim toq betda boşlanadimi; bob betida yuqori yozuv körinadimi; bob raqami va bezagi.
\item \textbf{Böş bet qoidasi} \beta\ --- bob oldidan qolgan böş betda yuqori va pastki yozuvlar bölmaydi.
\item \textbf{Fon elementlari} \beta\ --- ustaga rasm, rangli maydon yoki çiziq qöyiş; şu usta işlagan har bir betda çiqadi, juft betlarda kerak bölsa közgu kabi aks etadi; bob oldidagi böş betlarda çiqmaydi.
\item \textbf{Bölimga xos usta} \beta\ --- muqaddima, asosiy matn, ilovalar uçun turli ustalar; \texttt{\textbackslash sahifausta\{id\}} bilan almaşadi.
\item \textbf{Butun kitobga} \beta\ --- ustani bir marta özgartirsangiz, u işlatilgan hamma bet yangilanadi.
\end{itemize}

\rasmjoy{1.2cm}{Ekran: «Ustalar» paneli --- ikki betli yoyilma (juft/toq), yuqori yozuv kataklari: çapda «bob nomi», öngda «bet raqami»; röyxatda «Asosiy», «Bob boşlanişi», «Ilova» ustalari.}{Sahifa ustalari}

% ---------------------------------------------------------------------
\section{Uslublar}\label{sec:sd-uslub}
\begin{itemize}
\item \textbf{Xatboşi uslublari} \beta\ --- Sarlavha 1--3, Asosiy matn, Şeʼr qatori, Iqtibos, Izoh, Rasm osti yozuvi va öz uslublaringiz.
\item \textbf{Belgi uslublari} \beta\ --- urğu, atama, xorijiy söz, sahifa havolasi va öz uslublaringiz.
\item \textbf{Qöllaş} \beta\ --- matnni belgilab uslubni tanlaysiz; uslub LaTeX buyruği bölib kodga yoziladi.
\item \textbf{Bir joyda özgartiriş} \beta\ --- uslubni özgartirsangiz, kitobdagi hamma joy yangilanadi, çunki uslub \texttt{sahifa.sty} da bir marta taʼriflangan.
\item \textbf{Asosiy uslub} \beta\ --- uslub boşqasidan meros oladi («Iqtibos» = «Asosiy matn» + kursiv + surilma).
\end{itemize}

\begin{center}\footnotesize
\begin{tabularx}{\linewidth}{>{\bfseries}p{30mm}>{\raggedright\arraybackslash}p{58mm}X}
\toprule
Uslub & LaTeX'da & Izoh\\
\midrule
Sarlavha 1, 2, 3 & \texttt{\textbackslash chapter}, \texttt{\textbackslash section}, \texttt{\textbackslash subsection} & körinişi titlesec bilan; mundarija, «Faqat şu bob» va SyncTeX özgarmaydi\\
Asosiy matn & hujjatning odatiy şrifti & \texttt{\textbackslash setmainfont}, \texttt{\textbackslash fontsize}\\
Şeʼr qatori & \texttt{\textbackslash begin\{uslub\}\{sher\}} & qatorlar \texttt{\textbackslash\textbackslash} bilan; uzun qator surilib davom etadi\\
Iqtibos, Izoh & \texttt{\textbackslash begin\{uslub\}\{iqtibos\}} & ixtiyoriy nomli uslub ham şu yöl bilan\\
Rasm osti yozuvi & caption paketi sozlamasi & \texttt{\textbackslash rasmramka[izoh=...]}\\
Belgi uslubi & \texttt{\textbackslash matnuslubi\{urgu\}\{...\}} & \\
\bottomrule
\end{tabularx}
\end{center}

% ---------------------------------------------------------------------
\section{Tipografiya paneli}\label{sec:sd-tipo}
\begin{center}\footnotesize
\begin{tabularx}{\linewidth}{>{\bfseries}p{34mm}X>{\raggedright\arraybackslash}p{40mm}}
\toprule
Sozlama & Nima qiladi & LaTeX'da\\
\midrule
Şrift & yuklangan TTF/OTF yoki örnatilgan şrift, har bir uslubga alohida & \texttt{\textbackslash newfontfamily} (fontspec)\\
Ölçam va qator oralişi & punkt bilan & \texttt{\textbackslash fontsize\{..\}\{..\}}\\
Harf oralişi (tracking) & uslub uçun & fontspec \texttt{LetterSpace}\\
Tayanç çiziq töri & vizual rejimda tör çiziqlari körinadi (\beta); joylarni avtomatik qator oralişiga karrali qiliş \reja & \ref{sec:sd-xavf}-bölimdagi çeklov bilan\\
Böğin köçiriş & yoqiş/öçiriş, eng kam harf soni; özbekça qoidalar (lotin, kirill, Ö Ğ Ş Ç) & \texttt{\textbackslash uzbekhyphenmins}, \texttt{\textbackslash hyphenpenalty}\\
Yolğiz qatorlar & bet boşida yoki oxirida bitta qator qolmasin & \texttt{\textbackslash widowpenalty}, \texttt{\textbackslash clubpenalty}; LuaLaTeX'da lua-widow-control\\
Bosh harf (drop cap) & 2--4 qator balandlikda & lettrine\\
Bob bezaklari & bob sarlavhasi ostida va bob oxirida \texttt{* * *}, qisqa çiziq yoki gul (\beta) & \texttt{\textbackslash bobbezagi}, pifont\\
Mikrotipografiya & çetlarni tekislaş, harflarni ozgina keng/tor qiliş & microtype (XeLaTeX'da faqat çetni tekislaş)\\
\bottomrule
\end{tabularx}
\end{center}

% ---------------------------------------------------------------------
\section{Rasm ramkalari va matn qutilari}\label{sec:sd-ramka}
\begin{itemize}
\item \textbf{Taşlab qöyiş} \reja\ --- rasmni sahifaga taşlaysiz; \texttt{\textbackslash rasmramka} yoziladi, rasm \texttt{rasmlar/} papkasiga tuşadi (hozirgi rasm qöyiş qoidalari bilan).
\item \textbf{Kesiş} \reja\ --- ramka çetlarini Alt bilan tortib kesiladi; graphicx \texttt{trim} va \texttt{clip} ga yoziladi, asl fayl özgarmaydi.
\item \textbf{Matn bilan oralaş} \beta\ --- rasm çap yoki öngda, matn uni aylanib ötadi (wrapfig).
\item \textbf{Sahifaga mahkamlaş} \beta\ --- rasm yoki matn qutisi betning aniq joyida turadi (textpos, eso-pic).
\item \textbf{Töliq bet va çetgaça} \beta\ --- rasm butun betni egallaydi va kesiş çizigidan 3 mm taşqariga çiqadi (bleed).
\item \textbf{Rasm osti yozuvi} \beta\ --- ramka ostida yoki yonida; uslubi «Rasm osti yozuvi».
\item \textbf{Rasmli kitob maketi} \beta\ --- töliq betli rasm (\texttt{joy=fon}) + ustida matn qutisi (\texttt{matnquti}); ikki betli yoyilmaga bitta rasm (\texttt{joy=yoyilma}).
\end{itemize}

\textbf{Rasm ramkasi rejimlari}
\begin{center}\footnotesize
\begin{tabularx}{\linewidth}{>{\bfseries}p{28mm}X>{\raggedright\arraybackslash}p{36mm}}
\toprule
Rejim & Qanday turadi & Paket\\
\midrule
\texttt{joy=oqim} & matn içida, xatboşilar orasida & graphicx\\
\texttt{joy=oralash} & çap yoki öngda, matn aylanib ötadi & wrapfig\\
\texttt{joy=sahifa} & betning aniq koordinatasida & textpos, eso-pic\\
\texttt{joy=tola} & butun bet, çetgaça; bet raqami va yuqori yozuv yöq & eso-pic\\
\texttt{joy=yoyilma} & ikki betga bitta rasm (çap yarmi juft betda, öng yarmi toq betda) & eso-pic, graphicx \texttt{trim}\\
\bottomrule
\end{tabularx}
\end{center}

\rasmjoy{1.2cm}{Ekran: bolalar kitobi yoyilmasi --- ikki betni egallagan rasm, kesiş çizigidan taşqariga çiqqan qizil bleed hoşiyasi, rasm ustida yarim şaffof fonli matn qutisi va uning tutqiçlari.}{Töliq betli rasm va matn qutisi}

% ---------------------------------------------------------------------
\section{Kitob turlari uçun şablonlar}\label{sec:sd-shablon}
Har bir şablon özbekça (lotin, kirill, Ö Ğ Ş Ç) va xatosiz kompilyatsiya qilinişi şart. \reja
\begin{center}\footnotesize
\begin{tabularx}{\linewidth}{>{\bfseries}p{28mm}X>{\raggedright\arraybackslash}p{30mm}}
\toprule
Şablon & Nima bor & Kompilyator\\
\midrule
Roman, nasr \beta & bir ustun, bob boşida bosh harf, bob oxirida bezak, toq betda boşlanadigan bob & XeLaTeX, LuaLaTeX\\
Şeʼriyat \beta & şeʼr qatori uslubi, raqamsiz betlar, şeʼr markazda, uzun qator surilib davom etadi & XeLaTeX, LuaLaTeX\\
Bolalar rasmli kitobi \beta & töliq betli rasmlar, katta şrift, rasm ustida matn qutisi, kvadrat format & XeLaTeX, LuaLaTeX\\
Luğat \beta & ikki ustun, qalin bosh söz, betning yuqorisida birinçi va oxirgi bosh söz, harf sarlavhalari & XeLaTeX, LuaLaTeX\\
Diniy matn \beta & arabça matn öngdan çapga, özbekça tarjima yonma-yon (paracol), oyat raqami & LuaLaTeX, XeLaTeX\\
Oşpazlik kitobi \beta & «Masalliqlar» va «Tayyorlaş» ikki ustunda, taom rasmi & XeLaTeX, LuaLaTeX\\
Iş daftari \beta & yoziş çiziqlari, javob uçun joy, katak qoğoz (mavjud dizayn paketi bilan) & XeLaTeX, LuaLaTeX\\
Fotoalbom, katalog \beta & rasm töri (2, 3, 4 rasm; \texttt{rasmtori}, \texttt{\textbackslash torrasm}), rasm osti yozuvi, mahsulot kartoçkasi & XeLaTeX, LuaLaTeX\\
\bottomrule
\end{tabularx}
\end{center}

% ---------------------------------------------------------------------
\section{Muqova dizayneri}\label{sec:sd-muqova}
\begin{itemize}
\item \textbf{Üç qism} \beta\ --- old muqova, orqa muqova va kitob bel qismi (koreşok) bitta yoyilmada.
\item \textbf{Bel kengligi} \beta\ --- bet soni va qoğoz qalinligidan hisoblanadi: \(\text{kenglik} = \frac{\text{betlar}}{2} \times \text{bir varaq qalinligi}\) + muqova qöşimçasi.
\item \textbf{Bleed va kesiş çiziqlari} \beta\ --- har tomonda bleed (masalan 3 mm) va kesiş belgilari.
\item \textbf{CMYK va ISBN} \beta\ --- ranglar CMYK; ISBN ştrix-kodi «Bosmaxona» modulidagi EAN-13 generatori bilan.
\item \textbf{Alohida PDF} \beta\ --- muqova \texttt{muqova.tex} dan alohida kompilyatsiya qilinadi va alohida yuklab olinadi.
\end{itemize}

\textbf{Qoğoz qalinligi jadvali} \beta\ --- gazeta (45, 49), ofset (60--120), bör qoğoz yaltiroq va xira (90--170), kitob qoğozi (60--80 g/m²) uçun odatiy qalinliklar; ular taxminiy (zavodga qarab $\pm$10\,\%), şuning uçun bosmaxona bergan aniq qiymat alohida kiritiladi. Muqovalaş: yelimlangan yumşoq, skoba, qattiq (karton 2,5 mm, 3 mm çiqib turadi, 7 mm bukilma, 15 mm qayirma --- odatiy qiymatlar). \tekshir{qattiq muqova ölçamlarini bosmaxona bilan solişriş}

\rasmjoy{1.2cm}{Ekran: muqova yoyilmasi --- çapda orqa muqova (annotatsiya, ştrix-kod), örtada 14 mm bel (kitob nomi tik yozilgan), öngda old muqova; atrofida bleed va kesiş belgilari; yonida «Bet soni: 256, qoğoz: 80 g/m²» maydonlari.}{Muqova dizayneri}

% ---------------------------------------------------------------------
\section{Preflight: bosmadan oldingi tekşiruv}\label{sec:sd-preflight}
\tugma{Preflight} kompilyatsiya qilingan PDF ni tekşiradi va röyxat körsatadi; tuzatib bölsa, yonida tugma. \beta
\begin{center}\footnotesize
\begin{tabularx}{\linewidth}{>{\bfseries}p{36mm}X>{\raggedright\arraybackslash}p{40mm}}
\toprule
Tekşiruv & Qanday aniqlanadi & Bir bosişda tuzatiş\\
\midrule
Şriftlar joylanganmi & \texttt{pdffonts} & yöq (XeLaTeX/LuaLaTeX şriftlarni odatda joylaydi)\\
Rasm sifati (dpi) & \texttt{pdfimages -list}: joylangan ölçamdagi piksel ziçligi & yöq; rasm nomi va beti körsatiladi\\
Rasm rangi (RGB/CMYK) & \texttt{pdfimages -list} & yöq; bosmaxonaga ögiriş kerakligi aytiladi\\
Bleed & log da çetgaça rasm (\texttt{tola}, \texttt{fon}, \texttt{yoyilma}) bor, lekin PDF beti kesilgan betdan katta emas & yöq; «Bosmaxona»da kesiş belgilarini yoqiş aytiladi\\
Bet soni 16 yoki 32 ga karrali & log dagi bet soni & oxiriga kerakli böş betlar qöşiş\\
Yolğiz qatorlar va qisqa betlar & yolğiz qatorlarni jarimalar oldini oladi; bet oxirida 1--2 qator qolgan betlar \texttt{pdftotext -layout} qatorlari bilan topiladi & yöq; \texttt{\textbackslash enlargethispage} maslahati\\
Sahifadan çiqqan qatorlar & log dagi «Overfull \textbackslash hbox» & yöq; qatorga ötiş\\
Böş betlar & matni ham, rasmi ham yöq bet (\texttt{\textbackslash bosbet} dan taşqari) & yöq; betga ötiş\\
\bottomrule
\end{tabularx}
\end{center}
Hammasi işlaydi: şriftlar, rasm sifati va rangi, bet soni (böş betlarni qöşiş tugmasi bilan), sahifadan çiqqan qatorlar, böş betlar, qisqa betlar va bleed.

% ---------------------------------------------------------------------
\section{Import}\label{sec:sd-import}
\begin{itemize}
\item \textbf{Word (.docx)} \beta\ --- matn, rasmlar va formulalar hozirgi «Wordʼdan import» kabi; qöşimça ravişda Word xatboşi uslublari (Heading 1, Quote, öz uslublaringiz) Oversells uslublariga aylanadi.
\item \textbf{InDesign (.idml)} \beta\ --- ramkalar sahifadagi tartibida öqiladi; matn, xatboşi va belgi uslublari (ölçam, qator oralişi, tekislaş, surilmalar, oraliqlar, qalin va kursiv, rang, bosh harf, meros), sarlavhalar, izohlar, sahifa ölçami va birinçi ustaning hoşiyalari ötadi. Ötmaydi: ramkalarning aniq joyi, jadval (matni xatboşi bölib qoladi), boğlangan rasmlar va şriftlar --- ular röyxati \texttt{main.tex} boşida.
\item \textbf{PDF} \beta\ --- faqat matn (\texttt{pdftotext -layout} sandbox içida): xatboşilar surilma va böş qatordan topiladi, takrorlanadigan sarlavha va bet raqamlari olib taşlanadi, satr oxiridagi bölinişlar tiklanadi («bir-biri» kabi juft sözlar saqlanadi), sarlavhalar taxminiy (\texttt{\textbackslash section*}). Uslublar, rasmlar va jadvallar ötmaydi; skanerlangan PDF da matn yöq.
\item \textbf{Uslublar moslaşuvi} \beta\ --- «Uslublar» bölimida istalgan uslubni boşqasiga qöşasiz (masalan, Word'dan kelgan «Quote» → «Iqtibos»): loyihadagi barça \texttt{.tex} fayllarda u işlatilgan joylar özgaradi, uslub \texttt{sahifa.sty} dan öçiriladi, unga asoslangan uslublar yangisiga asoslanadi.
\end{itemize}

% ---------------------------------------------------------------------
\section{AI dizayn}\label{sec:sd-ai}
Faqat server egasi AI ni yoqqan bölsa işlaydi; yoqilmagan bölsa, tugmalar körinmaydi va heç narsa taşqariga yuborilmaydi. \beta
\begin{itemize}
\item \textbf{Kitobdan dizayn} \reja\ --- mavjud kitobning 2--3 betini (PDF yoki rasm) yuklaysiz; uslublar, sahifa ölçami, hoşiyalar va ustalar taklif qilinadi. Hozirgi «Kitobdan dizayn» imkoniyatining kengaytmasi.
\item \textbf{Tavsifdan dizayn} \beta\ --- «bolalar kitobi, iliq ranglar, katta şrift» kabi tavsif yozasiz; dizayn taklif qilinadi.
\item \textbf{Taklif --- faqat taklif} \beta\ --- natija körinadi, qabul qilsangizgina \texttt{sahifa.sty} ga yoziladi; «Tarix» da oldingi holat saqlanadi.
\end{itemize}
\begin{center}\footnotesize
\begin{tabularx}{\linewidth}{>{\bfseries}p{40mm}p{40mm}X}
\toprule
Imkoniyat & Qayerda & Taşqariga nima ketadi\\
\midrule
AI: kitobdan sahifa dizayni & taşqi AI serveri (Anthropic) & siz yuklagan betlar (PDF yoki rasm), istaklaringiz, şriftlar röyxati\\
AI: tavsifdan dizayn & taşqi AI serveri & siz yozgan tavsif va şriftlar röyxati\\
\bottomrule
\end{tabularx}
\end{center}
Kitob matni, loyiha fayllari va boşqa betlar yuborilmaydi.
 bilan
%  qöşiladi. Preambuladagi \tugma, \maslahat, \tayyor, \beta, \reja,
%  \tekshir va \rasmjoy işlatiladi. Hamma imkoniyat hozir \reja.
% =====================================================================

\section{Sahifa dizayni: umumiy körinişi}\label{sec:sd}
\begin{itemize}
\item \textbf{Nima} \reja\ --- kitob içki sahifalarini siçqonça bilan joylaş: hoşiyalar, ustunlar, rasm ramkalari, matn qutilari, sahifa ustalari va uslublar. Har bir özgariş LaTeX kodiga yoziladi.
\item \textbf{Qanday kitoblar} \beta\ --- roman va nasr, şeʼriyat, bolalar rasmli kitobi, luğat, diniy matn (arabça, öngdan çapga), oşpazlik kitobi, iş daftari, fotoalbom, katalog.
\item \textbf{Natija} \reja\ --- tahrirlanadigan LaTeX va bosmaga tayyor PDF. Yopiq yoki ikkilik format yöq: barça sozlamalar loyihadagi \texttt{.sty} va \texttt{.tex} fayllarda turadi.
\item \textbf{Özbek tili} \reja\ --- yangi alifbo (Ö Ğ Ş Ç) uçun imlo, özbekça böğin köçiriş va kirill ↔ lotin ötkaziş dizayn qilingan sahifalarda ham işlaydi.
\item \textbf{Kompilyator} \beta\ --- XeLaTeX va LuaLaTeX; arabça (öngdan çapga) şablon ikkalasida ham işlaydi, LuaLaTeX tavsiya etiladi (\ref{sec:sd-shablon}-bölim).
\end{itemize}

\rasmjoy{1.2cm}{Ekran: çapda kod, öngda kompilyatsiya qilingan sahifa; sahifa ustida kök hoşiya çiziqlari, ustun çegaralari va rasm ramkasining tutqiçlari; yuqorida «Sahifa», «Ustalar», «Uslublar», «Tipografiya», «Preflight» yorliqlari.}{Sahifa dizayni rejimi}

% ---------------------------------------------------------------------
\section{Körinadigan sahifa rejimi}\label{sec:sd-vizual}
\begin{itemize}
\item \textbf{\tugma{Sahifa dizayni}} \reja\ --- PDF panelini tahrir rejimiga ötkazadi: kod çapda, sahifa öngda qoladi.
\item \textbf{Hoşiya va ustunlar} \beta\ --- hoşiya çiziqlari va ustun çegaralari siçqonça bilan suriladi; qiymat millimetrda yonida körinadi.
\item \textbf{Ramkalar} \beta\ --- rasm ramkasi va matn qutisi suriladi, kattalaştiriladi, kesiladi.
\item \textbf{Koddan sahifaga} \beta\ --- kodni özgartirsangiz, kompilyatsiyadan keyin sahifa va tutqiçlar yangilanadi.
\item \textbf{Sahifadan kodga} \beta\ --- tutqiçni qöyib yuborganingizda faqat şu ramkaning kalitlari (\texttt{x=}, \texttt{y=}, \texttt{w=}\ldots) özgaradi; qolgan kod harfma-harf joyida qoladi.
\item \textbf{Siz yozgan kod} \beta\ --- vizual tahrirlovçiga tegişli bölmagan kod sahifada qulf belgisi bilan körinadi; ustiga bossangiz, kodning öşa qatori oçiladi.
\item \textbf{Birgalikda} \beta\ --- özgarişlar boşqa tahrirlar kabi jonli sinxron qilinadi va «Tarix» ga tuşadi.
\end{itemize}

\textbf{Vizual tahrirlovçiga tegişli LaTeX qurilmalari}
\begin{center}\footnotesize
\begin{tabularx}{\linewidth}{>{\raggedright\arraybackslash}p{50mm}X}
\toprule
Qurilma & Nimani boşqaradi\\
\midrule
\texttt{sahifa.sty} (butunlay) & sahifa ölçami, hoşiyalar, ustunlar, ustalar, uslublar, şriftlar; boşida \texttt{\% oversells-sahifa:} izohida model\\
\texttt{\textbackslash sahifausta\{id\}} & şu joydan boşlab qaysi usta işlaşi\\
\texttt{\textbackslash begin\{uslub\}\{id\}} & xatboşi uslubi\\
\texttt{\textbackslash matnuslubi\{id\}\{...\}} & belgi uslubi\\
\texttt{\textbackslash rasmramka[...]\{fayl\}} & rasm: oqimda, matn bilan oralaş, sahifada joylaşgan yoki töliq sahifa\\
\texttt{\textbackslash begin\{matnquti\}[...]} & sahifaga mahkamlangan matn qutisi\\
\texttt{\textbackslash begin\{ustunlar\}[...]\{n\}} & bir qismdagi ustunlar soni\\
\texttt{\textbackslash bosbet} & ataylab böş qoldirilgan bet\\
\bottomrule
\end{tabularx}
\end{center}
Boşqa hamma narsa (sizning buyruqlaringiz, \texttt{tikzpicture}, jadvallar) faqat kod orqali tahrirlanadi.

% ---------------------------------------------------------------------
\section{Sahifa ustalari}\label{sec:sd-usta}
\begin{itemize}
\item \textbf{Juft va toq betlar} \beta\ --- har biriga alohida yuqori va pastki yozuvlar: bob nomi, bölim nomi, kitob nomi, bet raqami yoki matn.
\item \textbf{Bet raqami} \beta\ --- joyi (yuqori/past, içki/taşqi çet, örta), şakli (1, i, I), boşlaniş raqami.
\item \textbf{Bob boşlanişi} \beta\ --- bob doim toq betda boşlanadimi; bob betida yuqori yozuv körinadimi; bob raqami va bezagi.
\item \textbf{Böş bet qoidasi} \beta\ --- bob oldidan qolgan böş betda yuqori va pastki yozuvlar bölmaydi.
\item \textbf{Fon elementlari} \beta\ --- ustaga rasm, rangli maydon yoki çiziq qöyiş; şu usta işlagan har bir betda çiqadi, juft betlarda kerak bölsa közgu kabi aks etadi; bob oldidagi böş betlarda çiqmaydi.
\item \textbf{Bölimga xos usta} \beta\ --- muqaddima, asosiy matn, ilovalar uçun turli ustalar; \texttt{\textbackslash sahifausta\{id\}} bilan almaşadi.
\item \textbf{Butun kitobga} \beta\ --- ustani bir marta özgartirsangiz, u işlatilgan hamma bet yangilanadi.
\end{itemize}

\rasmjoy{1.2cm}{Ekran: «Ustalar» paneli --- ikki betli yoyilma (juft/toq), yuqori yozuv kataklari: çapda «bob nomi», öngda «bet raqami»; röyxatda «Asosiy», «Bob boşlanişi», «Ilova» ustalari.}{Sahifa ustalari}

% ---------------------------------------------------------------------
\section{Uslublar}\label{sec:sd-uslub}
\begin{itemize}
\item \textbf{Xatboşi uslublari} \beta\ --- Sarlavha 1--3, Asosiy matn, Şeʼr qatori, Iqtibos, Izoh, Rasm osti yozuvi va öz uslublaringiz.
\item \textbf{Belgi uslublari} \beta\ --- urğu, atama, xorijiy söz, sahifa havolasi va öz uslublaringiz.
\item \textbf{Qöllaş} \beta\ --- matnni belgilab uslubni tanlaysiz; uslub LaTeX buyruği bölib kodga yoziladi.
\item \textbf{Bir joyda özgartiriş} \beta\ --- uslubni özgartirsangiz, kitobdagi hamma joy yangilanadi, çunki uslub \texttt{sahifa.sty} da bir marta taʼriflangan.
\item \textbf{Asosiy uslub} \beta\ --- uslub boşqasidan meros oladi («Iqtibos» = «Asosiy matn» + kursiv + surilma).
\end{itemize}

\begin{center}\footnotesize
\begin{tabularx}{\linewidth}{>{\bfseries}p{30mm}>{\raggedright\arraybackslash}p{58mm}X}
\toprule
Uslub & LaTeX'da & Izoh\\
\midrule
Sarlavha 1, 2, 3 & \texttt{\textbackslash chapter}, \texttt{\textbackslash section}, \texttt{\textbackslash subsection} & körinişi titlesec bilan; mundarija, «Faqat şu bob» va SyncTeX özgarmaydi\\
Asosiy matn & hujjatning odatiy şrifti & \texttt{\textbackslash setmainfont}, \texttt{\textbackslash fontsize}\\
Şeʼr qatori & \texttt{\textbackslash begin\{uslub\}\{sher\}} & qatorlar \texttt{\textbackslash\textbackslash} bilan; uzun qator surilib davom etadi\\
Iqtibos, Izoh & \texttt{\textbackslash begin\{uslub\}\{iqtibos\}} & ixtiyoriy nomli uslub ham şu yöl bilan\\
Rasm osti yozuvi & caption paketi sozlamasi & \texttt{\textbackslash rasmramka[izoh=...]}\\
Belgi uslubi & \texttt{\textbackslash matnuslubi\{urgu\}\{...\}} & \\
\bottomrule
\end{tabularx}
\end{center}

% ---------------------------------------------------------------------
\section{Tipografiya paneli}\label{sec:sd-tipo}
\begin{center}\footnotesize
\begin{tabularx}{\linewidth}{>{\bfseries}p{34mm}X>{\raggedright\arraybackslash}p{40mm}}
\toprule
Sozlama & Nima qiladi & LaTeX'da\\
\midrule
Şrift & yuklangan TTF/OTF yoki örnatilgan şrift, har bir uslubga alohida & \texttt{\textbackslash newfontfamily} (fontspec)\\
Ölçam va qator oralişi & punkt bilan & \texttt{\textbackslash fontsize\{..\}\{..\}}\\
Harf oralişi (tracking) & uslub uçun & fontspec \texttt{LetterSpace}\\
Tayanç çiziq töri & vizual rejimda tör çiziqlari körinadi (\beta); joylarni avtomatik qator oralişiga karrali qiliş \reja & \ref{sec:sd-xavf}-bölimdagi çeklov bilan\\
Böğin köçiriş & yoqiş/öçiriş, eng kam harf soni; özbekça qoidalar (lotin, kirill, Ö Ğ Ş Ç) & \texttt{\textbackslash uzbekhyphenmins}, \texttt{\textbackslash hyphenpenalty}\\
Yolğiz qatorlar & bet boşida yoki oxirida bitta qator qolmasin & \texttt{\textbackslash widowpenalty}, \texttt{\textbackslash clubpenalty}; LuaLaTeX'da lua-widow-control\\
Bosh harf (drop cap) & 2--4 qator balandlikda & lettrine\\
Bob bezaklari & bob sarlavhasi ostida va bob oxirida \texttt{* * *}, qisqa çiziq yoki gul (\beta) & \texttt{\textbackslash bobbezagi}, pifont\\
Mikrotipografiya & çetlarni tekislaş, harflarni ozgina keng/tor qiliş & microtype (XeLaTeX'da faqat çetni tekislaş)\\
\bottomrule
\end{tabularx}
\end{center}

% ---------------------------------------------------------------------
\section{Rasm ramkalari va matn qutilari}\label{sec:sd-ramka}
\begin{itemize}
\item \textbf{Taşlab qöyiş} \reja\ --- rasmni sahifaga taşlaysiz; \texttt{\textbackslash rasmramka} yoziladi, rasm \texttt{rasmlar/} papkasiga tuşadi (hozirgi rasm qöyiş qoidalari bilan).
\item \textbf{Kesiş} \reja\ --- ramka çetlarini Alt bilan tortib kesiladi; graphicx \texttt{trim} va \texttt{clip} ga yoziladi, asl fayl özgarmaydi.
\item \textbf{Matn bilan oralaş} \beta\ --- rasm çap yoki öngda, matn uni aylanib ötadi (wrapfig).
\item \textbf{Sahifaga mahkamlaş} \beta\ --- rasm yoki matn qutisi betning aniq joyida turadi (textpos, eso-pic).
\item \textbf{Töliq bet va çetgaça} \beta\ --- rasm butun betni egallaydi va kesiş çizigidan 3 mm taşqariga çiqadi (bleed).
\item \textbf{Rasm osti yozuvi} \beta\ --- ramka ostida yoki yonida; uslubi «Rasm osti yozuvi».
\item \textbf{Rasmli kitob maketi} \beta\ --- töliq betli rasm (\texttt{joy=fon}) + ustida matn qutisi (\texttt{matnquti}); ikki betli yoyilmaga bitta rasm (\texttt{joy=yoyilma}).
\end{itemize}

\textbf{Rasm ramkasi rejimlari}
\begin{center}\footnotesize
\begin{tabularx}{\linewidth}{>{\bfseries}p{28mm}X>{\raggedright\arraybackslash}p{36mm}}
\toprule
Rejim & Qanday turadi & Paket\\
\midrule
\texttt{joy=oqim} & matn içida, xatboşilar orasida & graphicx\\
\texttt{joy=oralash} & çap yoki öngda, matn aylanib ötadi & wrapfig\\
\texttt{joy=sahifa} & betning aniq koordinatasida & textpos, eso-pic\\
\texttt{joy=tola} & butun bet, çetgaça; bet raqami va yuqori yozuv yöq & eso-pic\\
\texttt{joy=yoyilma} & ikki betga bitta rasm (çap yarmi juft betda, öng yarmi toq betda) & eso-pic, graphicx \texttt{trim}\\
\bottomrule
\end{tabularx}
\end{center}

\rasmjoy{1.2cm}{Ekran: bolalar kitobi yoyilmasi --- ikki betni egallagan rasm, kesiş çizigidan taşqariga çiqqan qizil bleed hoşiyasi, rasm ustida yarim şaffof fonli matn qutisi va uning tutqiçlari.}{Töliq betli rasm va matn qutisi}

% ---------------------------------------------------------------------
\section{Kitob turlari uçun şablonlar}\label{sec:sd-shablon}
Har bir şablon özbekça (lotin, kirill, Ö Ğ Ş Ç) va xatosiz kompilyatsiya qilinişi şart. \reja
\begin{center}\footnotesize
\begin{tabularx}{\linewidth}{>{\bfseries}p{28mm}X>{\raggedright\arraybackslash}p{30mm}}
\toprule
Şablon & Nima bor & Kompilyator\\
\midrule
Roman, nasr \beta & bir ustun, bob boşida bosh harf, bob oxirida bezak, toq betda boşlanadigan bob & XeLaTeX, LuaLaTeX\\
Şeʼriyat \beta & şeʼr qatori uslubi, raqamsiz betlar, şeʼr markazda, uzun qator surilib davom etadi & XeLaTeX, LuaLaTeX\\
Bolalar rasmli kitobi \beta & töliq betli rasmlar, katta şrift, rasm ustida matn qutisi, kvadrat format & XeLaTeX, LuaLaTeX\\
Luğat \beta & ikki ustun, qalin bosh söz, betning yuqorisida birinçi va oxirgi bosh söz, harf sarlavhalari & XeLaTeX, LuaLaTeX\\
Diniy matn \beta & arabça matn öngdan çapga, özbekça tarjima yonma-yon (paracol), oyat raqami & LuaLaTeX, XeLaTeX\\
Oşpazlik kitobi \beta & «Masalliqlar» va «Tayyorlaş» ikki ustunda, taom rasmi & XeLaTeX, LuaLaTeX\\
Iş daftari \beta & yoziş çiziqlari, javob uçun joy, katak qoğoz (mavjud dizayn paketi bilan) & XeLaTeX, LuaLaTeX\\
Fotoalbom, katalog \beta & rasm töri (2, 3, 4 rasm; \texttt{rasmtori}, \texttt{\textbackslash torrasm}), rasm osti yozuvi, mahsulot kartoçkasi & XeLaTeX, LuaLaTeX\\
\bottomrule
\end{tabularx}
\end{center}

% ---------------------------------------------------------------------
\section{Muqova dizayneri}\label{sec:sd-muqova}
\begin{itemize}
\item \textbf{Üç qism} \beta\ --- old muqova, orqa muqova va kitob bel qismi (koreşok) bitta yoyilmada.
\item \textbf{Bel kengligi} \beta\ --- bet soni va qoğoz qalinligidan hisoblanadi: \(\text{kenglik} = \frac{\text{betlar}}{2} \times \text{bir varaq qalinligi}\) + muqova qöşimçasi.
\item \textbf{Bleed va kesiş çiziqlari} \beta\ --- har tomonda bleed (masalan 3 mm) va kesiş belgilari.
\item \textbf{CMYK va ISBN} \beta\ --- ranglar CMYK; ISBN ştrix-kodi «Bosmaxona» modulidagi EAN-13 generatori bilan.
\item \textbf{Alohida PDF} \beta\ --- muqova \texttt{muqova.tex} dan alohida kompilyatsiya qilinadi va alohida yuklab olinadi.
\end{itemize}

\textbf{Qoğoz qalinligi jadvali} \beta\ --- gazeta (45, 49), ofset (60--120), bör qoğoz yaltiroq va xira (90--170), kitob qoğozi (60--80 g/m²) uçun odatiy qalinliklar; ular taxminiy (zavodga qarab $\pm$10\,\%), şuning uçun bosmaxona bergan aniq qiymat alohida kiritiladi. Muqovalaş: yelimlangan yumşoq, skoba, qattiq (karton 2,5 mm, 3 mm çiqib turadi, 7 mm bukilma, 15 mm qayirma --- odatiy qiymatlar). \tekshir{qattiq muqova ölçamlarini bosmaxona bilan solişriş}

\rasmjoy{1.2cm}{Ekran: muqova yoyilmasi --- çapda orqa muqova (annotatsiya, ştrix-kod), örtada 14 mm bel (kitob nomi tik yozilgan), öngda old muqova; atrofida bleed va kesiş belgilari; yonida «Bet soni: 256, qoğoz: 80 g/m²» maydonlari.}{Muqova dizayneri}

% ---------------------------------------------------------------------
\section{Preflight: bosmadan oldingi tekşiruv}\label{sec:sd-preflight}
\tugma{Preflight} kompilyatsiya qilingan PDF ni tekşiradi va röyxat körsatadi; tuzatib bölsa, yonida tugma. \beta
\begin{center}\footnotesize
\begin{tabularx}{\linewidth}{>{\bfseries}p{36mm}X>{\raggedright\arraybackslash}p{40mm}}
\toprule
Tekşiruv & Qanday aniqlanadi & Bir bosişda tuzatiş\\
\midrule
Şriftlar joylanganmi & \texttt{pdffonts} & yöq (XeLaTeX/LuaLaTeX şriftlarni odatda joylaydi)\\
Rasm sifati (dpi) & \texttt{pdfimages -list}: joylangan ölçamdagi piksel ziçligi & yöq; rasm nomi va beti körsatiladi\\
Rasm rangi (RGB/CMYK) & \texttt{pdfimages -list} & yöq; bosmaxonaga ögiriş kerakligi aytiladi\\
Bleed & log da çetgaça rasm (\texttt{tola}, \texttt{fon}, \texttt{yoyilma}) bor, lekin PDF beti kesilgan betdan katta emas & yöq; «Bosmaxona»da kesiş belgilarini yoqiş aytiladi\\
Bet soni 16 yoki 32 ga karrali & log dagi bet soni & oxiriga kerakli böş betlar qöşiş\\
Yolğiz qatorlar va qisqa betlar & yolğiz qatorlarni jarimalar oldini oladi; bet oxirida 1--2 qator qolgan betlar \texttt{pdftotext -layout} qatorlari bilan topiladi & yöq; \texttt{\textbackslash enlargethispage} maslahati\\
Sahifadan çiqqan qatorlar & log dagi «Overfull \textbackslash hbox» & yöq; qatorga ötiş\\
Böş betlar & matni ham, rasmi ham yöq bet (\texttt{\textbackslash bosbet} dan taşqari) & yöq; betga ötiş\\
\bottomrule
\end{tabularx}
\end{center}
Hammasi işlaydi: şriftlar, rasm sifati va rangi, bet soni (böş betlarni qöşiş tugmasi bilan), sahifadan çiqqan qatorlar, böş betlar, qisqa betlar va bleed.

% ---------------------------------------------------------------------
\section{Import}\label{sec:sd-import}
\begin{itemize}
\item \textbf{Word (.docx)} \beta\ --- matn, rasmlar va formulalar hozirgi «Wordʼdan import» kabi; qöşimça ravişda Word xatboşi uslublari (Heading 1, Quote, öz uslublaringiz) Oversells uslublariga aylanadi.
\item \textbf{InDesign (.idml)} \beta\ --- ramkalar sahifadagi tartibida öqiladi; matn, xatboşi va belgi uslublari (ölçam, qator oralişi, tekislaş, surilmalar, oraliqlar, qalin va kursiv, rang, bosh harf, meros), sarlavhalar, izohlar, sahifa ölçami va birinçi ustaning hoşiyalari ötadi. Ötmaydi: ramkalarning aniq joyi, jadval (matni xatboşi bölib qoladi), boğlangan rasmlar va şriftlar --- ular röyxati \texttt{main.tex} boşida.
\item \textbf{PDF} \beta\ --- faqat matn (\texttt{pdftotext -layout} sandbox içida): xatboşilar surilma va böş qatordan topiladi, takrorlanadigan sarlavha va bet raqamlari olib taşlanadi, satr oxiridagi bölinişlar tiklanadi («bir-biri» kabi juft sözlar saqlanadi), sarlavhalar taxminiy (\texttt{\textbackslash section*}). Uslublar, rasmlar va jadvallar ötmaydi; skanerlangan PDF da matn yöq.
\item \textbf{Uslublar moslaşuvi} \beta\ --- «Uslublar» bölimida istalgan uslubni boşqasiga qöşasiz (masalan, Word'dan kelgan «Quote» → «Iqtibos»): loyihadagi barça \texttt{.tex} fayllarda u işlatilgan joylar özgaradi, uslub \texttt{sahifa.sty} dan öçiriladi, unga asoslangan uslublar yangisiga asoslanadi.
\end{itemize}

% ---------------------------------------------------------------------
\section{AI dizayn}\label{sec:sd-ai}
Faqat server egasi AI ni yoqqan bölsa işlaydi; yoqilmagan bölsa, tugmalar körinmaydi va heç narsa taşqariga yuborilmaydi. \beta
\begin{itemize}
\item \textbf{Kitobdan dizayn} \reja\ --- mavjud kitobning 2--3 betini (PDF yoki rasm) yuklaysiz; uslublar, sahifa ölçami, hoşiyalar va ustalar taklif qilinadi. Hozirgi «Kitobdan dizayn» imkoniyatining kengaytmasi.
\item \textbf{Tavsifdan dizayn} \beta\ --- «bolalar kitobi, iliq ranglar, katta şrift» kabi tavsif yozasiz; dizayn taklif qilinadi.
\item \textbf{Taklif --- faqat taklif} \beta\ --- natija körinadi, qabul qilsangizgina \texttt{sahifa.sty} ga yoziladi; «Tarix» da oldingi holat saqlanadi.
\end{itemize}
\begin{center}\footnotesize
\begin{tabularx}{\linewidth}{>{\bfseries}p{40mm}p{40mm}X}
\toprule
Imkoniyat & Qayerda & Taşqariga nima ketadi\\
\midrule
AI: kitobdan sahifa dizayni & taşqi AI serveri (Anthropic) & siz yuklagan betlar (PDF yoki rasm), istaklaringiz, şriftlar röyxati\\
AI: tavsifdan dizayn & taşqi AI serveri & siz yozgan tavsif va şriftlar röyxati\\
\bottomrule
\end{tabularx}
\end{center}
Kitob matni, loyiha fayllari va boşqa betlar yuborilmaydi.


\part*{B. Texnik loyiha}
% =====================================================================
\section{Ma’lumotlar modeli}\label{sec:tl-model}
\begin{itemize}
\item \textbf{Bitta manba} --- model loyihadagi \texttt{sahifa.sty} faylining boşida, \texttt{\% oversells-sahifa:} bilan boşlanadigan izoh qatorlarida JSON bölib turadi. Çiziş taxtasidagi \texttt{\% oversells-drawing:} usuli bilan bir xil.
\item \textbf{Fayl özi LaTeX} --- izohdan keyingi qism modeldan hosil qilinadi; Oversells bölmasa ham fayl odatiy LaTeX paketi sifatida kompilyatsiya qilinadi.
\item \textbf{Ramkalar matnda} --- rasm ramkasi va matn qutisining joyi \texttt{\textbackslash rasmramka[...]} kalitlarida, öşa bob faylining özida turadi; alohida maʼlumotlar bazasi yöq.
\item \textbf{Tarix va birgalikda işlaş} --- hamma narsa matn fayl bölgani uçun «Tarix», versiyalarni soliştiriş va jonli sinxron özgarişsiz işlaydi.
\item \textbf{Qölda yozilgan qöşimçalar} --- foydalanuvçi özi yozadigan sozlamalar \texttt{sahifa-qoshimcha.tex} da; Oversells uni heç qaçon qayta yozmaydi.
\end{itemize}

\textbf{Model tarkibi} (MVP)
\begin{center}\footnotesize
\begin{tabularx}{\linewidth}{>{\bfseries}p{34mm}X}
\toprule
Maydon & Mazmuni\\
\midrule
\texttt{version} & model versiyasi (1)\\
\texttt{paragraphStyles[]} & \texttt{id} (lotin kiçik harf va raqam), nom, \texttt{basedOn}, şrift, ölçam, qator oralişi, harf oralişi, tekislaş, surilma, oldingi/keyingi joy, kursiv, qalin, rang, böğin köçiriş, keyingisi bilan birga, bosh harf\\
\texttt{characterStyles[]} & \texttt{id}, nom, kursiv, qalin, kiçik bosh harflar, rang, şrift, harf oralişi\\
\texttt{masters[]} & \texttt{id}, nom; toq va juft bet uçun yuqori va pastki yozuv, har birida çap, örta, öng katak (böş, bet raqami, bob, bölim, kitob nomi, matn); bet raqami şakli\\
\texttt{defaultMaster} & kitob boşidan işlaydigan usta\\
\texttt{chapterMaster} & bob boşlanadigan betning ustasi\\
\texttt{chapterOpensRight} & bob doim toq betda boşlanadi\\
\texttt{blankPagesEmpty} & bob oldidagi böş betda yozuvlar yöq\\
\bottomrule
\end{tabularx}
\end{center}

\begin{Verbatim}
% oversells-sahifa: {"version":1,"paragraphStyles":[{"id":"iqtibos",
% oversells-sahifa: "name":"Iqtibos","italic":true,"indent":10}],...}
\ProvidesPackage{sahifa}[Oversells sahifa dizayni]
...
\end{Verbatim}

% =====================================================================
\section{LaTeX hosil qiliş}\label{sec:tl-latex}
\begin{center}\footnotesize
\begin{tabularx}{\linewidth}{>{\bfseries}p{38mm}>{\raggedright\arraybackslash}p{40mm}X}
\toprule
Qism & Paket & Qanday\\
\midrule
Xatboşi uslubi & xparse (LaTeX yadrosi) & \texttt{\textbackslash begin\{uslub\}\{id\}}: \texttt{\textbackslash csname} orqali şu uslub sozlamalari qöllanadi; meros \texttt{basedOn} bilan\\
Belgi uslubi & --- & \texttt{\textbackslash matnuslubi\{id\}\{matn\}}\\
Sarlavhalar & titlesec & Sarlavha 1--3 \texttt{\textbackslash chapter}/\texttt{\textbackslash section}/\texttt{\textbackslash subsection} körinişini özgartiradi, nomlarini emas\\
Ustalar & fancyhdr & har bir usta --- \texttt{\textbackslash fancypagestyle\{sahifa@id\}}; \texttt{\textbackslash sahifausta\{id\}} uni yoqadi\\
Bob boşlanişi & titlesec, fancyhdr & bob betida \texttt{chapterMaster}; \texttt{chapterOpensRight} da \texttt{\textbackslash cleardoublepage}\\
Böş betlar & --- & \texttt{\textbackslash cleardoublepage} böş betni \texttt{empty} uslubida qoldiradi\\
Şriftlar & fontspec & \texttt{\textbackslash newfontfamily}; yuklangan fayl uçun \texttt{Path=}\\
Tipografiya & microtype, lettrine, lua-widow-control & 2-bosqiç\\
Ramkalar & graphicx, wrapfig, textpos, eso-pic & 2-bosqiç\\
Arabça (RTL) & babel \texttt{bidi=basic}, paracol & faqat LuaLaTeX; XeLaTeX uçun kerakli \texttt{bidi} paketi kompilyatsiya obrazida yöq\\
Muqova & tikz, crop, «Bosmaxona» moduli & alohida \texttt{muqova.tex}\\
\bottomrule
\end{tabularx}
\end{center}
\begin{itemize}
\item \textbf{Ulaniş} --- asosiy faylga \texttt{\textbackslash usepackage\{sahifa\}} qöşiladi: \texttt{\textbackslash begin\{document\}} dan oldin, \texttt{dizayn} paketi bölsa undan keyin.
\item \textbf{Tekşirilgan} --- kompilyatsiya obrazida bor: textpos, wrapfig, lettrine, paracol, flowfram, eso-pic, grid-system, pdfx, lua-widow-control, nowidow, crop, microtype; yöq: bidi, memoir, verse, reledmac.
\item \textbf{Tekşirilgan} --- özbekça böğin köçiriş yangi harflarda işlaydi: boş-lanğiç-la-ri-miz-dan, öz-gar-ti-riş-la-ri-miz-dan.
\item \textbf{Tekşirilgan} --- LuaLaTeX, babel \texttt{bidi=basic}, Noto Naskh Arabic va paracol bilan arabça matn öngdan çapga terildi. Noto Naskh Arabic da oddiy qavs \texttt{( )} yöq: oyat raqami uçun bezakli qavslar (U+FD3E, U+FD3F) yoki boşqa şrift kerak.
\end{itemize}

% =====================================================================
\section{Ikki tomonlama sinxron}\label{sec:tl-sync}
\begin{enumerate}
\item \textbf{Koddan sahifaga.} Kompilyatsiyada \texttt{sahifa.sty} har bir ramka uçun \texttt{zref-savepos} bilan bet raqami va koordinatalarni yozadi; server ularni JSON qilib beradi; brauzer PDF ustiga tutqiçlarni şu joyda çizadi. Hoşiya va ustun çiziqlari kompilyatsiyani kutmay modeldan çiziladi.
\item \textbf{Sahifadan kodga.} Tutqiç qöyib yuborilganda yangi qiymat millimetrda hisoblanadi; kodda faqat şu ramkaning \texttt{[...]} qismi almaştiriladi (id böyiça topiladi). Özgariş boşqa tahrirlar kabi Yjs orqali saqlanadi.
\item \textbf{Kutiş vaqti.} Ramka darhol yangi joyda xira körinadi; haqiqiy sahifa «Avto» kompilyatsiyadan keyin keladi. Matn qayta oqişini InDesign kabi bir zumda körsatib bölmaydi.
\item \textbf{Uslub qöllaş.} Kodda: belgilangan matn uslub buyruği bilan örab olinadi. PDF da: SyncTeX matnni qator aniqligida topadi --- xatboşi uslubi uçun yetarli; belgi uslubi uçun belgilangan söz öşa qatorlarda qidiriladi, bir neça joyda uçrasa yoki buyruq içida bölsa, kod paneli taklif qilinadi.
\item \textbf{Buzilgan qurilma.} Vizual tahrirlovçi tanimaydigan şaklga keltirilgan \texttt{\textbackslash rasmramka} «qölda yozilgan» deb belgilanadi va tutqiç olmaydi.
\end{enumerate}

% =====================================================================
\section{Özbek tili dizayn içida}\label{sec:tl-til}
\begin{itemize}
\item \textbf{Ötkazgiç} --- \texttt{latex-text.ts} da \texttt{uslub} muhiti argumenti va \texttt{\textbackslash matnuslubi}, \texttt{\textbackslash sahifausta} birinçi argumenti kalit sifatida belgilanadi: kirill ↔ lotin va Ö Ğ Ş Ç ötkaziş \texttt{id} ni buzmaydi, matnni ötkazadi.
\item \textbf{Imlo} --- uslub buyruqlari içidagi matn odatiy matn kabi tekşiriladi.
\item \textbf{Usta matni} --- yuqori yozuvdagi erkin matn modelda saqlanadi; yozuvni ötkazişda \texttt{sahifa.sty} qayta hosil qilinadi \tekshir{MVP da qölda; 2-bosqiçda avtomatik}.
\end{itemize}

% =====================================================================
\section{500 betdan katta kitoblar}\label{sec:tl-tezlik}
\begin{itemize}
\item \textbf{«Faqat şu bob»} --- vizual rejimdagi har bir özgariş faqat kursor turgan bobni kompilyatsiya qiladi (\texttt{main-bob.tex}).
\item \textbf{Kesh} --- mundarija va havolalar oldingi kompilyatsiyadan olinadi (\texttt{.oversells-cache.tar}).
\item \textbf{Rasm nusxalari} --- tahrir paytida katta rasmlar örniga kiçik nusxalar (masalan, 150 dpi) işlatiladi; preflight va yakuniy PDF asl rasmlar bilan \tekshir{tezlikka taʼsiri ölçansin}.
\item \textbf{Koordinatalar} --- faqat kompilyatsiya qilingan bob uçun yangilanadi.
\item \textbf{Model kiçik} --- uslublar va ustalar bir neça kilobayt; saqlaş va ötkaziş kitob hajmiga boğliq emas.
\end{itemize}

% =====================================================================
\section{Mavjud modullar bilan boğlaniş}\label{sec:tl-integratsiya}
\begin{itemize}
\item \textbf{Dizayn paketi} (\texttt{dizayn.sty}) --- matematika darsliklari uçun qutilar, formulalar, çizmalar. \texttt{sahifa.sty} undan keyin yuklanadi va faqat foydalanuvçi özgartirgan narsani (sarlavha körinişi, ustalar, şriftlar) qayta belgilaydi; dizayn ranglari uslublarda işlatilişi mumkin.
\item \textbf{Bosmaxona moduli} (\texttt{bosma.sty}) --- kitob formati, kesiş çiziqlari, CMYK, ISBN. Sahifa ölçamining yagona manbai şu modul; \texttt{sahifa.sty} faqat hoşiyalar va ustunlarni beradi. Muqova dizayneri uning EAN-13 generatorini işlatadi.
\item \textbf{AI} --- hozirgi «Kitobdan dizayn» endpointi kengaytiriladi; javob \texttt{sahifa.sty} modeli şaklida.
\end{itemize}

\textbf{Kod joylaşuvi (MVP)}
\begin{center}\footnotesize
\begin{tabularx}{\linewidth}{>{\raggedright\arraybackslash}p{78mm}X}
\toprule
Fayl & Vazifa\\
\midrule
\texttt{packages/shared/src/page-design.ts} & model turlari, tekşiriş, \texttt{sahifa.sty} hosil qiliş va öqiş, uslub qöllaş\\
\texttt{apps/server/src/page-design/} & \texttt{GET}/\texttt{PUT /api/projects/:id/page-design}\\
\texttt{apps/web/src/components/PageDesignDialog.tsx} & «Uslublar» va «Ustalar» oynasi\\
\texttt{PAGE\_DESIGN\_ENABLED} & funksiya bayroği; \texttt{false} bölsa, endpointlar 404 va tugma yöq\\
\bottomrule
\end{tabularx}
\end{center}

\part*{C. Foydalanuvçi yöllari}
% =====================================================================
\section{Öqituvçi: iş daftari}
\begin{enumerate}
\item «Iş daftari» şablonidan loyiha oçadi.
\item \tugma{Sahifa dizayni} → «Uslublar»: «Topşiriq» uslubiga 13 pt va qalin raqam beradi.
\item «Ustalar»: toq betda öngda bet raqami, çapda «2-sinf. Matematika»; juft betda aksinça.
\item Topşiriq matnini belgilab «Topşiriq» uslubini qöllaydi.
\item Javob joyi kerak bölgan joyga \tugma{Çizma va formulalar} dan «yoziş çiziqlari»ni qöyadi.
\item Preflight: bet soni 16 ga karrali emas → «Böş bet qöşiş» ni bosadi.
\item \tugma{Bosmaxona} dan A4 formatda PDF oladi.
\end{enumerate}

\section{Naşriyot: roman}
\begin{enumerate}
\item Muallifning \texttt{.docx} faylini «Wordʼdan import» bilan yuklaydi; Word uslublari jadvalda Oversells uslublariga moslanadi.
\item «Roman» şabloni ustalarini qöllaydi: bob toq betda, bob betida yuqori yozuv yöq, böş betlar toza.
\item Tipografiya: matn şrifti (yuklangan OTF), 10,5/14 pt, yolğiz qatorlar taqiqlangan, bob boşida 3 qatorli bosh harf.
\item \tugma{Bosmaxona}: 60×84/16, kesiş çiziqlari, CMYK.
\item Preflight: şriftlar joylangan, «Overfull» qatorlar röyxati bilan tuzatadi.
\item Muqova dizayneri: 320 bet va qoğoz qalinligidan bel kengligi hisoblanadi; muqova alohida PDF.
\end{enumerate}

\section{Bolalar kitobi muallifi: töliq betli rasmlar}
\begin{enumerate}
\item «Bolalar rasmli kitobi» şabloni (kvadrat format, katta şrift).
\item Rasmni sahifaga taşlaydi → rejim «töliq bet»: rasm bleedgaça kengayadi.
\item Rasm ustiga matn qutisi qöyadi, uni siçqonça bilan joylaydi; quti fonini yarim şaffof qiladi.
\item Ikki betli rasm uçun «yoyilma» rejimini tanlaydi.
\item Preflight: 300 dpi dan past rasmlar va RGB rasmlar röyxati.
\item Muqova va içki betlarni alohida PDF qilib bosmaxonaga beradi.
\end{enumerate}

\part*{D--F}
% =====================================================================
\section{Xavflar va çeklovlar}\label{sec:sd-xavf}
\begin{center}\footnotesize
\begin{tabularx}{\linewidth}{>{\bfseries\raggedright\arraybackslash}p{36mm}XX}
\toprule
Qiyin joy & Nima uçun & Yeçim yoki oçiq çeklov\\
\midrule
Matnning ramkadan ramkaga oqişi & LaTeX matnni bitta oqimda teradi; InDesign dagi boğlangan ramkalar yöq & oddiy kitoblarda kerak emas; murakkab maket uçun flowfram tekşiriladi \tekshir{flowfram ni sinab köriş}; aks holda çeklov deb aytiladi\\
Aniq tayanç çiziq töri & rasm, formula, sarlavha qator balandligini buzadi & joylar qator oralişiga karrali; LuaLaTeX da qöşimça tuzatiş \tekshir{}; 100\,\% aniqlik vaʼda qilinmaydi\\
RTL aralaş matn & XeLaTeX uçun bidi paketi obrazda yöq & arabça kitoblar faqat LuaLaTeX; bidi ni obrazga qöşiş alohida qaror\\
Bir zumda körsatiş & har özgariş kompilyatsiya talab qiladi & «Faqat şu bob», kesh, kiçik rasm nusxalari; ramka darhol xira körinadi\\
Belgi uslubini PDF da qöllaş & SyncTeX faqat qatorni biladi & noaniq bölsa kod paneli\\
Qölda yozilgan kod & foydalanuvçi kodini buzib qöyiş xavfi & faqat tanilgan qurilmalar özgaradi; qolgani qulf\\
IDML import & formati keng; pandoc öqimaydi & matn va uslublar; ramka joyi taxminiy\\
PDF dan import & PDF da tuzilma yöq & faqat matn\\
Muqova bel kengligi & qoğoz qalinligi turliça & naşriyot qiymat kiritadi; taxmin yozilmaydi\\
Şrift litsenziyasi & yuklangan şriftni PDF ga joylaş huquqi & foydalanuvçi javobgar; preflight ogohlantiradi \tekshir{matni}\\
\bottomrule
\end{tabularx}
\end{center}

% =====================================================================
\section{Yöl xaritasi}\label{sec:sd-reja}
\begin{center}\footnotesize
\begin{tabularx}{\linewidth}{>{\bfseries\raggedright\arraybackslash}p{40mm}XXX}
\toprule
Imkoniyat & MVP & 2-bosqiç & 3-bosqiç\\
\midrule
1. Vizual rejim & --- & hoşiya, ustun, ramka tutqiçlari; koordinatalar & belgi uslubini PDF da qöllaş\\
2. Ustalar & toq/juft yozuvlar, bet raqami, bob betida usta, böş bet qoidasi, bölimga xos usta & fon elementlari, vizual tahrir & ---\\
3. Uslublar & xatboşi va belgi uslublari, meros, koddan qöllaş & PDF dan qöllaş (xatboşi) & ---\\
4. Tipografiya & şrift, ölçam, qator va harf oralişi, böğin köçiriş, yolğiz qatorlar & bosh harf, bezaklar, microtype, tayanç töri & ---\\
5. Ramkalar & --- & oqim, oralaş, sahifa, töliq bet, kesiş & yoyilma\\
6. Şablonlar & --- & roman, şeʼriyat, iş daftari, bolalar kitobi & luğat, diniy (LuaLaTeX), oşpazlik, fotoalbom, katalog\\
7. Muqova & --- & --- & muqova dizayneri\\
8. Preflight & --- & şriftlar, dpi, bet soni, Overfull, böş betlar & bleed, CMYK, yolğiz qatorlar, tuzatiş tugmalari\\
9. Import & --- & Word uslublari & IDML, PDF matni\\
10. AI dizayn & --- & --- & kitobdan va tavsifdan\\
\bottomrule
\end{tabularx}
\end{center}
MVP \texttt{PAGE\_DESIGN\_ENABLED} bayroği ortida çiqadi va mavjud loyihalarni özgartirmaydi.

% =====================================================================
\section{Siz uçun savollar}\label{sec:sd-savol}
\begin{enumerate}
\item Arabça va RTL kitoblar uçun \texttt{bidi} paketini kompilyatsiya obraziga qöşamizmi (XeLaTeX uçun) yoki faqat LuaLaTeX bilan çeklanamizmi?
\item Qaysi kitob turlari birinçi navbatda kerak: roman, şeʼriyat, iş daftari, bolalar kitobi --- shu tartibmi?
\item Naşriyotlar qaysi qoğoz turlarini işlatadi va bir varaq qalinligi qança (muqova bel kengligi uçun)?
\item IDML import haqiqatda kerakmi? Qaysi naşriyot InDesign fayllarini beradi?
\item Bleed hajmi har doim 3 mm mi yoki bosmaxonaga qarab özgaradimi?
\item Tahrir paytida kiçik rasm nusxalari işlatilsa, preflight va yakuniy PDF asl rasmlar bilan bölişi yetarlimi?
\item Uslub nomlari interfeysda özbekça bölsin, \texttt{id} lar esa lotin kiçik harflarda (\texttt{iqtibos}, \texttt{sher}) --- shu qoida mos keladimi?
\item MVP qaysi foydalanuvçilarda sinaladi (bayroq qaysi serverda yoqiladi)?
\end{enumerate}

\end{document}
